\subsection{V)}
{\large\begin{align*}\displaystyle F(A, B, C, D) &= \sum(3, 7, 11, 13, 14, 15) \qquad x = 6; 2\end{align*}}

\subsubsection{Obtener la mínima expresión booleana}
\begin{center}
\begin{tabular}{c|cccc|c}
\rowcolor[HTML]{FFCCC9} 
\cellcolor[HTML]{9AFF99}$m$ & $A$ & $B$ & $C$ & $D$ & \cellcolor[HTML]{FFFFC7}$F(A,B,C,D)$ \\ \hline
0                           & 0   & 0   & 0   & 0   &                                    \\
1                           & 0   & 0   & 0   & 1   &                                    \\
2                           & 0   & 0   & 1   & 0   & $x$                                    \\
3                           & 0   & 0   & 1   & 1   & 1                                  \\
4                           & 0   & 1   & 0   & 0   &                                    \\
5                           & 0   & 1   & 0   & 1   &                                    \\
6                           & 0   & 1   & 1   & 0   & $x$                                   \\
7                           & 0   & 1   & 1   & 1   & 1                                  \\
8                           & 1   & 0   & 0   & 0   &                                    \\
9                           & 1   & 0   & 0   & 1   &                                    \\
10                          & 1   & 0   & 1   & 0   &                                    \\
11                          & 1   & 0   & 1   & 1   & 1                                  \\
12                          & 1   & 1   & 0   & 0   &                                    \\
13                          & 1   & 1   & 0   & 1   & 1                                  \\
14                          & 1   & 1   & 1   & 0   & 1                                  \\
15                          & 1   & 1   & 1   & 1   & 1                                 
\end{tabular}
\end{center}
\begin{center}
    \begin{karnaugh-map}[4][4][1][$C\qquad D$][$A\qquad B$]
        \minterms{3,7,11,13,14,15}
        \terms{2,6}{$x$}
        \implicant{13}{15}
        \implicant{3}{6}
        \implicant{3}{11}
        \implicant{7}{14}
    \end{karnaugh-map}
\end{center}

\begin{itemize}[nosep]
    \item Primer bloque ($m_3$, $m_2$, $m_7$ y $m_6$) $B$ y $D$ cambian por lo que salen y hay que negar $A$. 
    \item Segundo bloque ($m_7$, $m_6$, $m_15$ y $14$) $A$ y $D$ cambian por lo que salen.
    \item Tercer bloque ($m_13$ y $m_15$) cambia $C$.
    \item Cuarto bloque ($m_3$, $m_7$, $m_15$ y $m_11$) salen $A$ y $B$.
    \item La función simplificada queda:
\end{itemize}

{\large\begin{align*}F(A,B,C,D) &= \overline{A}\cdot C + B\cdot C + A\cdot B\cdot D + C\cdot D\end{align*}}
\noindent Se podría sacar factor común $C$:
{\large\begin{align*}F(A,B,C,D) &= C\cdot(\overline{A}+B+D)+A\cdot{}B\cdot{}D\end{align*}}

\subsubsection{Simulación}
\imagen[]{10cm}{./imagenes/ejercicio5Falstad.png}
\subsubsection{Descripción en verilog}
\begin{lstlisting}
module descripcionVerilogEjercicio5(
    input inputA,
    input inputB,
    input inputC,
    input inputD,
    output outputFunction
    );
	 assign outputFunction = inputC & (~inputA | inputB | inputD) | inputA & inputB & inputD;

endmodule
\end{lstlisting}
\subsubsection{RTL output}\imagen[Bloque RTL]{15cm}{./imagenes/rtlEjercicio5Verilog.png}\imagen[Implementación con compuertas]{15cm}{./imagenes/rtlCompuertasEjercicio5Verilog.png}\subsubsection{Simulación temporal}\noindent Código del módulo Verilog Test Fixture:\begin{lstlisting}
module verilogTestFixtureEjercicio5;

	// Inputs
	reg inputA;
	reg inputB;
	reg inputC;
	reg inputD;

	// Outputs
	wire outputFunction;

	// Instantiate the Unit Under Test (UUT)
	descripcionVerilogEjercicio5 uut (
		.inputA(inputA), 
		.inputB(inputB), 
		.inputC(inputC), 
		.inputD(inputD), 
		.outputFunction(outputFunction)
	);

	initial begin
		// Initialize Inputs
		inputA = 0;
		inputB = 0;
		inputC = 0;
		inputD = 0;

		// Wait 100 ns for global reset to finish
		#100;
      inputA = 0; inputB = 0; inputC = 0; inputD = 1;
		#100;
      inputA = 0; inputB = 0; inputC = 1; inputD = 0;
		#100;
      inputA = 0; inputB = 0; inputC = 1; inputD = 1;
		#100;
      inputA = 0; inputB = 1; inputC = 0; inputD = 0;
		#100;
      inputA = 0; inputB = 1; inputC = 0; inputD = 1;
		#100;
      inputA = 0; inputB = 1; inputC = 1; inputD = 0;
		#100;
      inputA = 0; inputB = 1; inputC = 1; inputD = 1;
		#100;
      inputA = 1; inputB = 0; inputC = 0; inputD = 0;
		#100;
      inputA = 1; inputB = 0; inputC = 0; inputD = 1;
		#100;
      inputA = 1; inputB = 0; inputC = 1; inputD = 0;
		#100;
      inputA = 1; inputB = 0; inputC = 1; inputD = 1;
		#100;
      inputA = 1; inputB = 1; inputC = 0; inputD = 0;
		#100;
	   inputA = 1; inputB = 1; inputC = 0; inputD = 1;
		#100;
      inputA = 1; inputB = 1; inputC = 1; inputD = 0;
		#100;
      inputA = 1; inputB = 1; inputC = 1; inputD = 1;

		// Add stimulus here

	end
      
endmodule\end{lstlisting}\imagen[Simulación temporal]{15cm}{./imagenes/simulacionTemporalEjercicio5.png}

